% !PRINTED none
% Abschnitt-opening page: no running page number visible in any of the three bands
% (top band is blank page-top whitespace; the title block sits lower on the page).
\begin{center}
\gesperrt{Vierzehnter Abschnitt.}

\textbf{Komposition der Formen und Ideale.}

\rule{0.2\linewidth}{0.4pt}

\S\,101. \gesperrt{Komposition in den Ordnungen.}
\end{center}

Es sei $\Delta$ die K\"orperdiskriminante und $D = Q^2\Delta$ die Diskriminante der Ordnung $[Q]$.

Ist $T = (a, b, c)$ eine primitive quadratische Form der Diskriminanten $D$ % [sic] Gedruckt ist Diskriminanten im Plural.
mit positivem, zu $Q$ teilerfremdem $a$, so k\"onnen wir daraus eine Basisform $\lambda$ in $O$ eines zu $Q$ teilerfremden Ideals $\mathfrak{a}$ in folgender Weise bestimmen:

Wir setzen
\weq{1}{\lambda = a t_1 + \frac{b + \sqrt{D}}{2} t_2,}
\weq{2}{N(\lambda) = a(a t_1^2 + b t_1 t_2 + c t_2^2).}
Bezeichnen wir das durch das Funktional $\lambda$ bestimmte Ideal, d. h. den gr\"o\ss ten gemeinschaftlichen Teiler von
\weq{3}{a, \frac{b + \sqrt{D}}{2}}
mit $\mathfrak{a}$, so folgt aus (2)
\weq{4}{N(\mathfrak{a}) = a.}
Hierin ist $a$ die kleinste durch $\mathfrak{a}$ teilbare nat\"urliche Zahl. Denn ist diese kleinste Zahl $a_1$, so mu\ss{} $a_1$ ein Teiler von $a$ sein, also etwa $a = a_1 a_2$. Hierin k\"onnen $a_1, a_2, b$ keinen gemeinsamen Teiler haben, denn sie k\"onnen erstens nicht alle drei gerade sein, weil sonst $a = 4a'$, $b = 2b'$, $c = c'$,
\[
D = 4(b'^2 - 4a'c)
\]
w\"are, und 4 m\"u\ss te in $Q$ aufgehen, entgegen der Annahme, da\ss{} $a$ zu $Q$ teilerfremd sei; und ebenso w\"urde ein ungerader gemeinschaftlicher Primteiler von $a_1, b, a_2$ in $Q$ aufgehen.
